%% kapa: exam=1 round=37 year=2026 period=2 form=A subject=accounting subject_ko=회계학 seq=18 number=58
%% source: Q-Net 감정평가사 기출문제 (cst00309) artlSeq=5258924 file=raw/1008/5258924/2026년 제37회 감정평가사 1차 2교시.pdf
%% conversion: kordoc 4.24.51 -> tools/gen1.py (auto); review=proofread-sonnet; flags=none
\begin{question}{58}
(주)감평은 20×1년 12월 31일 고객에게 1개월 이내에 반품가능한 조건으로 원가 ￦1,000,000의 상품을 ￦1,300,000에 현금 판매하였다. (주)감평은 위 거래금액의 2 \%가 반품될 것으로 신뢰성 있게 추정하였고 반품될 가능성은 높다. 상품이 반품될 경우 상품원가의 10 \%가 손상되며, 반품 관련 직접비용이 ￦5,000 발생할 것으로 추정하였다. 동 반품조건 판매 관련 회계처리가 (주)감평의 20×1년 당기순이익에 미치는 영향은?
\begin{choices}
  \item ￦195,000 증가
  \item ￦213,000 증가
  \item ￦285,000 증가
  \item ￦287,000 증가
  \item ￦295,000 증가
\end{choices}
\end{question}
\answer{4}
